% ==========================================================
% 1. PLANTEAMIENTO DEL PROBLEMA
% Esta es la sección más avanzada del documento: 1.2 ya está
% redactada con fuentes reales; el resto sigue en construcción.
% ==========================================================

\section{Planteamiento del problema}

\subsection{Antecedentes}
\estado{En construcción}
\notafase{Se redacta en la Fase 2 del curso.}

\subsection{Descripción de la situación problemática}

La problemática de esta investigación se relaciona con las condiciones en que los
estudiantes de Ingeniería de Sistemas desarrollan sus actividades académicas,
especialmente aquellas que requieren un uso prolongado del computador y otros
dispositivos electrónicos. Las extensas jornadas dedicadas al estudio, la
programación, la consulta de información y el desarrollo de proyectos pueden generar
molestias físicas y visuales, como fatiga ocular, sequedad o irritación de los ojos,
dolores o molestias en el cuello, hombros, espalda, brazos y muñecas, además de
cansancio físico y mental, particularmente cuando estas actividades se realizan
acompañadas de posturas inadecuadas, permanencia prolongada en una misma posición y
pausas insuficientes [1], [2]. En población universitaria se ha reportado una elevada
frecuencia de estas afectaciones; en particular, el síndrome visual informático
constituye una problemática frecuente entre estudiantes expuestos de manera
prolongada a dispositivos electrónicos [2], mientras que los trastornos
musculoesqueléticos presentan una alta frecuencia entre universitarios, con
afectaciones principalmente en regiones como el cuello, la zona lumbar y la espalda
[1]. Esta situación puede verse favorecida por la exposición continua a pantallas,
debido al esfuerzo visual sostenido y a la ausencia de periodos adecuados de
descanso; por la adopción de posturas inadecuadas, como mantener el cuello inclinado
hacia adelante, la espalda encorvada, los hombros elevados o las muñecas en
posiciones poco favorables; y por la permanencia durante largos periodos en una
misma posición. Asimismo, la falta de pausas activas y descansos visuales puede
reducir las oportunidades de recuperación física durante las jornadas académicas,
mientras que unas condiciones ergonómicas inadecuadas del puesto de estudio, como
mobiliario que no se ajusta correctamente, computadores portátiles ubicados a una
altura inconveniente o dispositivos de entrada utilizados en posiciones poco
favorables, pueden contribuir a la aparición de molestias musculoesqueléticas [1]. A
estas condiciones se suman factores relacionados con el entorno de estudio, como una
iluminación inadecuada, reflejos o una configuración poco apropiada de la pantalla,
que pueden incrementar el esfuerzo visual durante sesiones prolongadas frente a
dispositivos electrónicos [2]. Como consecuencia de estas condiciones, los
estudiantes pueden experimentar fatiga visual, manifestada mediante sensación de ojos
cansados, secos, irritados o visión borrosa; molestias o dolor en cuello y hombros
asociados con posturas estáticas y prolongadas; molestias en espalda, brazos, manos y
muñecas relacionadas con la permanencia prolongada frente al computador y con
posiciones poco adecuadas durante el uso del teclado y el mouse; así como cansancio
físico y mental, que puede afectar el bienestar y la concentración durante las
actividades académicas [1], [3]. De igual manera, la exposición prolongada a
dispositivos electrónicos puede relacionarse con dolores de cabeza y otras molestias
asociadas al uso intensivo de pantallas [3]. Aunque estas manifestaciones pueden
presentarse con frecuencia durante la formación universitaria, algunos estudiantes
pueden llegar a normalizarlas como parte de las exigencias de su carrera y no
reconocer oportunamente la importancia de adoptar medidas preventivas. En este
contexto, la realización de pausas activas y cambios periódicos de posición
constituye una estrategia que puede favorecer la recuperación durante las jornadas
académicas y contribuir al bienestar físico de los estudiantes [4]. Por esta razón,
resulta pertinente analizar la relación existente entre las condiciones de estudio,
los hábitos frente al computador y las molestias físicas y visuales que pueden
experimentar los estudiantes de Ingeniería de Sistemas, con el propósito de
identificar los principales factores asociados y contribuir al planteamiento de
medidas preventivas que favorezcan condiciones más adecuadas para el desarrollo de
sus actividades académicas.

\bigskip
\begin{center}
\begin{longtable}{c L{7cm} L{5cm}}
\caption{Relación entre síntomas y causas} \\
\toprule
\textbf{Ord.} & \textbf{Síntoma} & \textbf{Causas relacionadas} \\
\midrule
\endfirsthead
1 & Fatiga visual después de jornadas prolongadas frente a pantallas & 1, 2, 3, 4 \\
2 & Sensación de ojos secos, cansados o irritados & 1, 2, 3 \\
3 & Dolor o molestia en cuello y hombros & 1, 3, 4, 5 \\
4 & Dolor o molestia en espalda & 1, 3, 4, 5 \\
5 & Molestias en muñecas, brazos o manos & 1, 3, 4, 5, 6 \\
6 & Dolores de cabeza y disminución de la comodidad durante el estudio & 1, 2, 3 \\
\bottomrule
\end{longtable}
\end{center}

\bigskip
\begin{center}
\begin{longtable}{c L{7cm} L{5cm}}
\caption{Relación entre causas y síntomas} \\
\toprule
\textbf{Ord.} & \textbf{Causa} & \textbf{Síntomas relacionados} \\
\midrule
\endfirsthead
1 & Exposición prolongada a pantallas durante las jornadas de estudio & 1, 2, 3, 4, 5 \\
2 & Falta de pausas o descansos visuales durante el estudio & 1, 2 \\
3 & Permanecer durante largos periodos en una misma posición & 1, 2, 3, 4, 5 \\
4 & Adopción de posturas inadecuadas durante el uso del computador & 3, 4, 5, 6 \\
5 & Puesto de estudio con condiciones ergonómicas inadecuadas & 3, 4, 5, 6 \\
6 & Falta de pausas activas y cambios de posición & 3, 4, 5, 6 \\
\bottomrule
\end{longtable}
\end{center}

\subsection{Formulación del problema}
\estado{En construcción}
\notafase{Se redacta en la Fase 2 del curso.}

\subsection{Consecuencias de no investigar}
\estado{En construcción}
\notafase{Se redacta en la Fase 2 del curso.}

\subsection{Delimitación}
\estado{En construcción}
\notafase{Se redacta en la Fase 2 del curso.}
